\documentclass[a4paper]{article}
\usepackage[polish]{babel}
\usepackage[cp1250]{inputenc}
\usepackage[T1]{fontenc}
\usepackage[dvips]{graphicx}

\pagestyle{headings}
\textwidth      15.5cm
\oddsidemargin    .1cm
\evensidemargin   .1cm

\begin{document}

\thispagestyle{empty}

\begin{minipage}{5cm}
\includegraphics[scale=0.7]{salomon.eps}
\end{minipage}
\begin{minipage}{10cm}
\begin{center}
{\Large Salomon\\
System przetwarzania wiedzy\\}
\end{center}
\end{minipage}

\vspace*{0.5cm}

\hrulefill

\vspace*{1cm}


\begin{flushright}
{\Large \emph{Risk Analysis, ver 1.0} }
\end{flushright}



\begin{flushleft}
{ \Large Historia wersji }\[\]
\begin{tabular}{|p{2.5cm}||p{1.5cm}|p{6cm}|p{3.5cm}|}
  \hline
  Data & Wersja & Opis & Autor \\
  \hline \hline
  26.04.2005 & 1.0 & utworzenie & \emph{Leszek Grzanka}\\
  \hline
  \hline
\end{tabular}
\end{flushleft}

\newpage

\section{Ryzyko}

\subsection{Wdra�anie nowych cz�onk�w zespo�u}
Pewne op�nienia czasowe spowodowane by�y nieprzewidzianymi trudno�ciami z wdro�eniem narz�dzi u�ywanych przy pracy z projektem, takich jak \emph{Ant}, \emph{Subversion}. Z powodu braku wyspecyfikowania w dokumentacji projektowej, jakich konkretnie wersji narz�dzi trzeba u�ywa� (tak, aby dostosowa� je do u�ywanej \emph{Javy 1.5}), powsta�y problemy zwi�zane z kompilacj� i instalacj� platformy \emph{Salomon ver 0.9}.

Ryzyko nie by�o przewidziane.

Aby z�agodzi� skutki wyst�pienia tego ryzyka odby�o si� spotkanie z zespo�em developer�w w celu wyja�nienia ich w�tpliwo�ci.

\subsection{Ryzyko wydajno�ciowe}
Istotnym czynnikiem mog�cym wp�yn�� na warto�� ko�cowego produktu mo�e by� niedostosowanie platformy do rozmiaru przetwarzanych danych. Mo�liwe jest i� demonstracyjne wtyczki nie poradz� sobie z ekstrakcj� wiedzy z du�ej ilo�ci danych. Konieczne b�dzie przeprowadzenie test�w wydajno�ciowych.


\end{document}

